\documentclass[10pt,a4paper,twocolumn]{article}
\usepackage[revtexheadings=section,authoryear]{natureastro-custom}

\title{A Nature Astronomy-inspired manuscript in \LaTeX}
\author{First Author$^{1}$, Second Author$^{2}$ \& Third Author$^{1}$}
% \date{}  % hide date if not needed
% \runninghead{First Author et al.}

\begin{abstract}
  This unofficial template is intended for writing manuscripts, notes and polished preprints in a visual style inspired by Nature Astronomy. The current setup combines Palatino-style serif text and matching mathematics with compact two-column text, REVTeX-inspired section headings, author--year citations and concise captions. A title, author and abstract block sits above both columns. Journal production metadata such as received, accepted or publication dates, DOI information and volume details are omitted.
\end{abstract}

\affiliations{%
  $^{1}$School of Physics, University of Example, Melbourne, VIC, Australia.\\
$^{2}$Institute for Astronomy, Example University, Example City, Country.}
\correspondence{corresponding.author@example.edu}

\begin{document}
\maketitle

A manuscript can begin directly with the scientific narrative rather than with a heading labelled ``Introduction''. The opening paragraphs should establish the context, identify the unresolved problem and state the contribution succinctly. The selected \texttt{authoryear} option enables author--year citations through \texttt{natbib}; use \verb|\citep| for parenthetical citations \citep{DongMelatos2024,DongMelatos2025,DongEtAl2026} and \verb|\citet| for textual citations, as in \citet{DongMelatos2024}. The reference list is generated from \texttt{examples.bib} with the \texttt{aasjournalv7} BibTeX style.

This file is deliberately a clean authoring template rather than a simulation of the journal's production workflow. Its purpose is to make drafts pleasant to read while keeping the source conventional and easy to adapt later.

The \texttt{authoryear} option enables author--year citations; omitting it gives numerical superscripts. Use \texttt{natureheadings} for the default heading style. To apply REVTeX styling through a chosen level, set \texttt{revtexheadings=} to \texttt{section}, \texttt{subsection} or \texttt{subsubsection}. The shorthand \texttt{revtexheadings} styles all three levels.

\section{Results}
\subsection{A compact two-column hierarchy}
The principal body is set in two columns. Sections have centred, bold, uppercase headings with Roman numerals. Subsections remain left-aligned and bold, and subsubsections left-aligned and italic, both without displayed numbers. Text and mathematics use matching Palatino-style fonts. A simple polynomial illustrates a single-column equation:
\begin{equation}
  f(x) = ax^2 + bx + c,
\end{equation}
where $a$, $b$ and $c$ are constant coefficients.

For an equation that needs both columns, the \texttt{widetext} environment
interrupts the two-column flow with a full-width display:
\begin{widetext}
  \begin{equation}
    \label{eq:wide}
    \begin{aligned}
      I &= \int_{x_0}^{x_1} f(x)\,\mathrm{d}x
      = \int_{x_0}^{x_1} \left(ax^2 + bx + c\right)\,\mathrm{d}x,\\
      &= \left[\frac{a}{3}x^3 + \frac{b}{2}x^2 + cx\right]_{x_0}^{x_1}
      = \frac{a}{3}\left(x_1^3-x_0^3\right)
      + \frac{b}{2}\left(x_1^2-x_0^2\right) + c\left(x_1-x_0\right).
    \end{aligned}
  \end{equation}
\end{widetext}
Here $I$ is the definite integral of $f(x)$ between $x_0$ and $x_1$.
Text resumes in two columns after equation~\ref{eq:wide}; the short rules
mark the transition into and out of the wide display.

\begin{figure}[t]
  \centering
  \fbox{\rule{0pt}{35mm}\rule{0.94\columnwidth}{0pt}}
  \caption{\textbf{Example single-column figure.} Replace the framed placeholder with \texttt{\textbackslash includegraphics}. Captions should explain symbols and uncertainties sufficiently for the figure to be interpreted without repeatedly returning to the main text.}
  \label{fig:example}
\end{figure}

Figure~\ref{fig:example} illustrates the default single-column treatment. For a figure spanning the whole page width, use the standard \texttt{figure*} environment.

\subsection{Tables}
Tables use compact text and simple horizontal rules provided by \texttt{booktabs}. The example lists illustrative coefficients for the polynomial above; use \texttt{table*} for a table spanning both columns.
\begin{table}[t]
  \caption{Illustrative polynomial coefficients}
  \centering
  \begin{tabular}{lccc}
    \toprule
    Example & $a$ & $b$ & $c$ \\
    \midrule
    A & 1 & 0 & 0 \\
    B & 1 & 2 & 1 \\
    C & 2 & $-1$ & 3 \\
    \bottomrule
  \end{tabular}
\end{table}

\section{Discussion}
Use the layout as a writing environment for a manuscript, internal draft, research note or preprint. The source remains ordinary \LaTeX{}, so the manuscript can later be transferred to a journal-specific template without restructuring the content.

\methods
Methods can continue in the same two-column format. For a real manuscript, place sufficient implementation detail here to reproduce the analysis, together with statistical conventions, software versions and selection criteria.

\dataavailability
State where the data underlying the study can be obtained, if this section is useful for the document.

\codeavailability
State where analysis code is archived, if applicable.

\acknowledgements
Acknowledge funding, facilities, collaborations and other contributions as needed.

\authorcontributions
F.A. conceived the study. F.A. and S.A. performed the analysis. All authors discussed the results and contributed to the manuscript.

\competinginterests
The authors declare no competing interests.

\bibliography{examples}% Produces the bibliography via BibTeX.
\bibliographystyle{aasjournalv7}

\makeendmatter
\end{document}
